\section{Background and Motivating Example}
\label{sec:background}

In this section, we first introduce the key concepts underlying our approach,
including the Python language features and the Hoare-Triple.
We then illustrate why type-error detection must track type-state changes
across these calls through a motivating example.

\subsection{Dynamic Typing and Duck Typing in Python}

Type-error detection first determines the types of the operands at an
operation and then checks whether the operation supports those
types~\cite{oh2024pyinder}. In Python, however, both steps are complicated by
its flexible type system. Because Python is dynamically typed, the types
that reach an operation can change during execution. Even when these types
are known, duck typing requires compatibility to be judged by the behavior
provided by the operands rather than by their class names alone. Moreover,
different operand types may invoke different implementations of an
operation and produce results of different types, which further affects the
types that reach subsequent operations.

More specifically, the \textbf{dynamic typing} feature of Python allows a variable,包括局部变量和全局变量,
to be bound to objects of different types during program execution~\cite{pythonExecutionModel,pythonDataModel}.
Consequently, the type of an operand at an operation depends on the
assignments executed before that operation. These changes are not xxx within a single function: a call may return a new value, update a global
variable, or modify an attribute of an object shared with its caller.
Therefore, determining the types at an operation requires tracking changes
both within a function and across function calls.

Duck typing further affects both compatibility checking and type tracking.
In Python, an object is compatible with an operation if it provides the
required behavior, regardless of its nominal class~\cite{pythonGlossary}.
Thus, objects of different classes may support the same operation, while the
same operator may invoke different implementations for different operand
types. A detector must identify the methods or protocols supported by the
operands to determine whether the operation is valid. At the same time, the
selected implementation may produce a different return type or state
change, which must be tracked when checking subsequent operations. Type
annotations can describe the possible operand types or their required
interfaces~\cite{pep484}. However, type-error detection must further identify
the correspondence between the input and output types of an operation. This
correspondence indicates whether the operation is valid for a given
combination of operand types and, if so, what type of result it produces.


\subsection{Hoare Triples as Function Summaries}
\label{sec:hoare-background}

Because a function call may change the types subsequently observed by its
caller, detecting type errors across calls requires propagating these
changes from the callee to the caller. A function summary can provide this
information without reproducing the callee's implementation at each call
site. To support type-error detection, the summary must describe not only
the conditions before the call but also the type changes established after
it. Hoare triples provide a standard way to express this before-and-after
relationship.

Hoare logic describes how executing a program changes its
state~\cite{hoare1969axiomatic,softwareFoundationsHoare}. A Hoare triple has
the form
\begin{equation}
\{P\}\ C\ \{Q\},
\label{eq:hoare-triple}
\end{equation}
where $P$ is a precondition, $C$ is a command, and $Q$ is a postcondition.
The triple states that if $P$ holds before $C$ executes and $C$ terminates
normally, then $Q$ holds afterward. For example, $P$ may state that
\texttt{msg.body} has type \texttt{str}; $C$ may encode and update the body;
and $Q$ may state that \texttt{msg.body} has type \texttt{bytes}. The triple
therefore describes both the state required before a command and the state
established after it.

A Hoare triple can therefore summarize a function without copying its body
to every call site. While its precondition records what must hold when the
function is called, its postcondition records the return value and any
changes visible to the caller. When a call occurs, the function's parameters
are matched with the caller's arguments, after which the postcondition can
be used to check the operations following the call. Moreover, Hoare triples
can be composed: if $\{P\}\ C_1\ \{R\}$ and $\{R\}\ C_2\ \{Q\}$ hold, the
state $R$ produced by $C_1$ becomes the state used to analyze $C_2$.

This form of summary is particularly useful for Python type-error detection.
Although a function signature describes parameter and return types, it may
not show that the function changes a global variable or an attribute of an
argument. A Hoare triple can capture such changes in its postcondition and
then apply them at the call site. As a result, the caller obtains the updated
type state needed to check its next operation. Hoare triples are therefore
suitable for summarizing and propagating type changes across Python function
calls.

\subsection{A Call That Changes an Operation's Validity}
\label{sec:motivating-example}
Figure~\ref{fig:motivating-example} illustrates why type-error detection
requires tracking the state established by a call. The \texttt{Message}
class declares \texttt{body: str | bytes}. Both encoders accept a string. The \texttt{encode} method of the
text encoder returns that string; the UTF-8 encoder returns bytes. The helper
\texttt{encode\_message} writes the encoder's result back to the same
message object and implicitly returns \texttt{None}.

\begin{figure}[!htbp]
\centering
\includegraphics[width=\linewidth]{motivatingExample.pdf}
\caption{A call changes the type state used by a subsequent operation.
Point 1 establishes a string body; point 2 relates the selected encoder to
the field update; point 3 has a bytes body, making concatenation with a
string fail. The lower-right panel reports the illustrated example's
outcomes for Mypy, Pyrefly, Pyright, and our approach, PyTT.}
\Description{The caller make_line raises ValueError when the message body
is not a string. Otherwise it calls encode_message with Utf8Encoder, then
passes msg.body plus a newline string to print. The helper writes the
encoder result into msg.body. TextEncoder returns a string and Utf8Encoder
returns bytes. Numbered points mark the string state before the call,
the encoder-dependent update, and the bytes state after the call. The
result panel shows no diagnostic for Mypy, Pyrefly, and Pyright and a type
error detected by PyTT.}
\label{fig:motivating-example}
\end{figure}

\paragraph{From the guard to the update.}
In \texttt{make\_line}, a non-string body triggers \texttt{ValueError}.
Every execution reaching point~1 has therefore passed the guard with a
string body. The subsequent call supplies that message and a
\texttt{Utf8Encoder} instance to the helper. At point~2, the helper's two
encoder cases connect the input choice to different updates: the text
encoder establishes a string body and the UTF-8 encoder establishes bytes.
Both writes satisfy the field's union annotation. The relevant difference
is the state that each call establishes for subsequent operations.

\paragraph{From the update to the failure.}
Instantiating the UTF-8 case at the shown call binds the helper's message
to the caller's object. Consequently, point~3 contains a bytes body.
The string fact at point~1 concerns the old contents and does not constrain
the updated field. Evaluation of \texttt{msg.body + '\textbackslash n'}
then raises \texttt{TypeError} before \texttt{print} is invoked. With
\texttt{make\_line} as the analysis entry, that exception escapes the
function. Using \texttt{TextEncoder} at the same call would instead
establish a string body and permit the concatenation.

This reasoning distinguishes three representations. The helper's return
type, \texttt{None}, gives no description of the field update. A set of
possible field types, \texttt{str | bytes}, captures the possible contents.
A relation between the selected encoder, incoming body, and outgoing body
identifies the contents established in the actual calling context.
Preserving that relation lets the analyzer invalidate the old field fact
and check the addition using its new value.

The result panel in Figure~\ref{fig:motivating-example} identifies our
approach as \emph{PyTT}. For the illustrated example, it records no
diagnostic from Mypy, Pyrefly, or Pyright and a type error detected by PyTT;
the associated tool versions and configurations are xxx. The mechanism
illustrated by this result is the propagation of the encoder-dependent
write to the operation that consumes the field. Section~\ref{sec:evaluation}
evaluates detection across the benchmark under the stated configurations.

